Let $m^*$ denote Lebesgue \emph{outer} measure on $\mathbb R$.
For each $x\in\mathbb R$, suppose we are given a set $A_x\subset\mathbb R$ such that:
- $A_x$ is bounded, and
- $m^*(A_x)<1$.

Call a set $X\subseteq\mathbb R$ \emph{independent} (for the map $x\mapsto A_x$) if for all distinct $x,y\in X$ we have $x\notin A_y$.

Question 1: Must there exist an \emph{infinite} independent set $X\subseteq\mathbb R$?

Question 2: If additionally each $A_x$ is \emph{closed} and has (Lebesgue) measure $m(A_x)<1$, must there exist an independent set of size $3$?
